\section{存在问题和设计机会}

尽管近年来三维生成技术与平台发展迅速，
相关系统在降低三维内容生产门槛方面取得了显著进展，
但在具体应用场景与功能链路方面仍然存在一定局限。
通过对当前主流三维生成平台与研究系统的分析可以发现，
现有工具在材质生成流程、模型适配能力以及交互方式等方面仍存在改进空间，
这也为本研究提供了明确的设计机会。

当前多数三维生成平台的核心流程主要围绕
“文本或图像输入—自动生成三维模型—结果预览与导出”展开，
其主要目标是实现从二维信息到完整三维模型的自动化生成。
然而，在实际的三维内容生产流程中，
设计人员往往已经拥有大量现成的三维模型资源，
这些模型通常只包含几何结构而缺乏高质量材质。
对于此类“白模”资产，
用户更希望能够直接为其生成或迁移合适的材质，
而不是重新生成一个新的三维模型。

目前大多数 AI 三维生成平台并未针对这一需求提供完整的解决方案。
例如，平台通常允许用户上传图片或输入文本生成新的三维模型，
但缺乏“参考模型 + 白模”这一双模型输入的材质迁移流程。
在这种情况下，
如果用户希望将某一参考材质应用到已有模型上，
往往需要借助多个不同工具完成。
常见的流程是：
用户先在图像生成或风格迁移平台中对参考图片进行生成或风格调整，
随后再将生成结果导入到三维生成平台或建模软件中，
再通过平台生成风格化的三维模型。
这一流程不仅操作繁琐，
而且需要在多个系统之间反复切换，
大大降低了整体工作效率。

基于上述问题，
本研究提出构建一种面向三维模型材质迁移任务的自动化系统。
该系统以“参考材质模型 + 白模输入”为核心，
通过分析参考模型的纹理信息与几何特征，
结合目标白模的结构信息生成适配的材质贴图，
从而实现从白模到完整材质模型的自动转换。

在系统设计层面，
本研究还将构建一个基于 Web 的交互平台，
将模型上传、材质生成与三维预览整合到统一的操作流程中。
用户只需上传参考模型与目标白模，
系统即可自动完成材质推断与生成，
并在三维视图中实时展示结果。
这种一体化的交互模式能够有效减少跨平台操作步骤，
提升三维资产生成流程的效率与可用性。